% Generated draft: requires mathematical and source review.
\documentclass[11pt]{article}
\usepackage[margin=25mm]{geometry}
\usepackage{fontspec}
\setmainfont{Latin Modern Roman}
\usepackage{amsmath,amssymb,mathtools,unicode-math}
\setmathfont{Latin Modern Math}
\usepackage{enumitem,needspace,xurl,hyperref}
\setlength{\parindent}{0pt}
\setlength{\parskip}{5pt}
\setlength{\emergencystretch}{3em}
\newcommand{\N}{\mathbb{N}}
\newcommand{\Z}{\mathbb{Z}}
\newcommand{\Q}{\mathbb{Q}}
\newcommand{\R}{\mathbb{R}}
\newcommand{\C}{\mathbb{C}}
\newcommand{\F}{\mathbb{F}}
\newcommand{\A}{\mathbb{A}}
\newcommand{\PP}{\mathbb{P}}
\newcommand{\E}{\mathbb{E}}
\newcommand{\eps}{\varepsilon}
\newcommand{\dd}{\,\mathrm{d}}
\begin{document}
\section*{UnsolvedMath ID: 700118}
\textbf{Source problem code:} catalogue entry
\subsection{Definitions and mathematical statement}
Let \(r\geq 1\), let \(v=v(x_1,\ldots,x_r)\) be a group word, and let \(G\) be an arbitrary group. Define its set of values and its verbal subgroup by
\[
V_v(G)=\{v(g_1,\ldots,g_r):g_1,\ldots,g_r\in G\},\qquad vG=\langle V_v(G)\rangle.
\]
The marginal subgroup \(v^*G\) consists of all \(a\in G\) such that, for every \((g_1,\ldots,g_r)\in G^r\) and every \(1\leq i\leq r\),
\[
\begin{aligned}
v(g_1,\ldots,g_i a,\ldots,g_r)&=v(g_1,\ldots,g_i,\ldots,g_r),\\
v(g_1,\ldots,a g_i,\ldots,g_r)&=v(g_1,\ldots,g_i,\ldots,g_r).
\end{aligned}
\]
\textbf{Part (b), remaining finiteness question.} Is it true that, for every group word \(v\) and every group \(G\),
\[
[G:v^*G]<\infty\quad\Longrightarrow\quad |vG|<\infty?
\]
The original conjecture attributed to P. Hall asserted more: if \([G:v^*G]=m<\infty\), then \(vG\) is finite and there exists an integer \(k\geq 0\) such that \(|vG|\mid m^k\). The supplied editors' comment reports counterexamples to the divisibility assertion, attributed to Yu. G. Kleiman, while retaining the finiteness question as open. These are supplied source assessments, not independently verified conclusions.

\textbf{Archived parts, retained for context.} The source also records the following conjectures attributed to P. Hall:
\begin{itemize}
\item[(a)] For every group word \(v\) and group \(G\), if \(V_v(G)\) is finite, then \(vG\) is finite.
\item[(c)] For every group word \(v\) and group \(G\) satisfying the maximum condition on subgroups, if \(vG\) is finite, then \([G:v^*G]<\infty\). Here the maximum condition means that every ascending chain of subgroups of \(G\) eventually stabilizes.
\end{itemize}
The supplied archive reports that (a) and (c) are refuted, citing S. V. Ivanov and I. S. Ashmanov with A. Yu. Olshanskii, respectively. They are historical conjectures rather than additional open targets in this draft. The source does not expand the phrase ``maximum condition''; its usual interpretation on all subgroups is used above.
\subsection{Short English statement}
For an arbitrary group word \(v\) and group \(G\), does finite index of the marginal subgroup \(v^*G\) force the verbal subgroup \(vG\) to be finite? The supplied source labels this remaining question open, reports that the stronger divisibility claim fails, and archives parts (a) and (c) as refuted. The overall registry status is \texttt{partially\_solved}; these assessments require human source and mathematical review.
\subsection{Sources}
\begin{itemize}
\item Supplied local UnsolvedMath registry record, permanent ID \(700118\), exact source problem code \texttt{catalogue entry}, title \emph{Kourovka Notebook Problem 2.45}. The record names Yu. I. Merzlyakov as proposer, with proposed year 1967; its individual conjectures are attributed to P. Hall.
\item \emph{The Kourovka Notebook (Unsolved Problems in Group Theory)}, 2nd issue (1967), Problem 2.45; supplied as updated in the 21st issue (2026), \texttt{arXiv:1401.0300v48}. The collection description names E. I. Khukhro and V. D. Mazurov as editors.
\item Supplied editors' comment for part (b): Yu. G. Kleiman, \emph{Trans. Moscow Math. Soc.}, 1983, no. 2, 63--110; cited for examples in which \(|vG|\) does not divide a power of \([G:v^*G]\).
\item Supplied Archive of Solved Problems entry for part (a): S. V. Ivanov, \emph{Soviet Math. (Izv. VUZ)}, \textbf{33}, no. 6 (1989), 59--70; cited as refuting (a).
\item Supplied Archive of Solved Problems entry for part (c): I. S. Ashmanov and A. Yu. Olshanskii, \emph{Soviet Math. (Izv. VUZ)}, \textbf{29}, no. 11 (1985), 65--82; cited as refuting (c).
\item The supplied status review is dated \texttt{2026-10-16}, gives canonical status \texttt{partially\_solved}, and records evidence as ``published.'' This is supplied dated curation, not an independently verified literature assessment. No source was visited for this draft.
\end{itemize}
\end{document}
